\section{Giới Thiệu}\label{sec:introduction}

Trong kỹ nghệ phần mềm hướng mô hình, yêu cầu và giải pháp của một hệ thống thường được đặc tả bằng
hai loại mô hình tách biệt ở giai đoạn phân tích: mô hình ý định (goal model) trả lời hệ thống muốn
đạt được điều gì và vì sao, còn mô hình tiến trình (process model) trả lời hệ thống làm gì, theo
trình tự nào, để đạt được điều đó. Hai loại mô hình này thường do hai vai trò khác nhau xây dựng,
requirement engineer và business analyst, dùng những ngôn ngữ mô hình hoá khác nhau tuỳ dự án, và
tiến hoá độc lập theo thời gian khi hệ thống thay đổi.

Tiến hoá độc lập kéo theo một khoảng cách khó tránh giữa mô hình ý định và mô hình tiến trình, và
khoảng cách đó không dừng ở vấn đề tài liệu: đây là nguyên nhân trực tiếp gây phát sinh lại chi phí
và lỗi hiện thực hoá trong các dự án phần mềm phức tạp, khi chỉ một phần nhỏ hệ thống thông tin thực
sự được phát triển đúng theo tiến trình nghiệp vụ mà chúng lẽ ra phải phục
vụ~\cite{alignment2019slr}. Bài toán đặt ra, do đó, rất cụ thể: kiểm tra một mô hình tiến trình đã
cho có thật sự hiện thực hoá đúng ý định đã đặc tả trong một mô hình ý định đã cho hay không, kể cả
khi việc phân công vai trò tổ chức, ai được đảm nhiệm vai trò nào và đồng thời đảm nhiệm những vai
trò nào, không được đặc tả trong bất kỳ mô hình nào của cả hai. Phần~\ref{sec:motivation} cụ thể hoá
đầy đủ vấn đề này bằng một ví dụ.

Bài toán kiểm tra quan hệ hiện thực hoá giữa mô hình ý định và mô hình tiến trình không mới, và ba
hướng tiếp cận đại diện, trên ba cặp ngôn ngữ khác nhau, minh hoạ rõ mức trưởng thành hiện tại,
khảo sát đầy đủ ở Phần~\ref{sec:related-work}. Ghose và Koliadis~\cite{ghose2007validation} dịch
goal model GRL và process model BPMN sang Description Logics rồi dùng một reasoner phân loại tính
nhất quán tự động, nhưng suy luận dừng ở một thời điểm, không biểu diễn được thứ tự thực thi.
Koliadis và Ghose~\cite{koliadis2006goalbpm} đối chiếu goal model KAOS với process model BPMN bằng
cách gán mỗi hoạt động một effect annotation rồi đối chiếu hiệu ứng tích luỹ dọc từng đường thực thi
với đặc tả thời gian của goal, nhưng toàn bộ quy trình thực hiện hoàn toàn thủ công.
Caballero-Villalobos và cộng sự~\cite{caballero2026alig} tiến xa nhất: hình thức hoá goal model
iStar và process model, workflow net hoặc DCR graph, thành hai hệ chuyển trạng thái, kiểm compliance
bằng duyệt product LTS kèm chứng minh hình thức về tính đúng đắn và đầy đủ, nhưng buộc mọi nhánh
thực thi phải dẫn tới trạng thái thoả mọi quality mới được xem là compliant, một ràng buộc mà chính
nhóm tác giả thừa nhận có thể tạo ra phán quyết vô lý về nghiệp vụ ở những nhánh vốn dĩ không nên
thoả goal ban đầu. Dù khác nhau về ngôn ngữ cụ thể, cả ba hướng tiếp cận đều chia sẻ một giới hạn
gốc: không goal model hay process model nào trong số đó có chỗ để biểu diễn ràng buộc vai trò tổ
chức, nên không hướng nào phát hiện được vi phạm chỉ lộ ra ở tầng phân công vai trò.

Khoảng trống đó gọi ra một ngôn ngữ đặc tả cấu trúc tổ chức nằm trong cùng miền phân tích với mô
hình ý định và mô hình tiến trình, và đó là đề xuất trung tâm của bài báo: ngôn ngữ ACL. Bài báo
đánh giá đề xuất trên một cặp ngôn ngữ cụ thể, mô hình ý định i*~\cite{yu1997towards} và mô hình
tiến trình nghiệp vụ BPMN~\cite{omg2013bpmn}, hai ngôn ngữ phổ biến bậc nhất cho từng vế và đã có
sẵn ngữ nghĩa đủ chi tiết để suy snapshot tự động, nhưng nguyên lý đề xuất không giới hạn ở riêng
cặp ngôn ngữ này. Về bản chất kỹ thuật kiểm tra, phương pháp đề xuất thuộc đúng nhóm ngữ nghĩa vận
hành theo vết thực thi của GoalBPM và của Caballero-Villalobos và cộng sự, đối chiếu trạng thái vận
hành của tiến trình theo thời gian với đặc tả goal, thay vì suy luận tĩnh tại một thời điểm như kỹ
thuật Description Logics. Điểm khác biệt nằm ở cách suy ra trạng thái cần đối chiếu: phương pháp kế
thừa khung khái niệm kịch bản--hành động--snapshot của Dang và cộng sự~\cite{dang2010jucs}, nhưng
suy snapshot tự động từ các ràng buộc OCL đã có sẵn trong mô hình, thay vì viết tay như GoalBPM hay
tự xây luật lan truyền riêng như Caballero-Villalobos và cộng sự. Ràng buộc cấu trúc tổ chức của ACL
được dịch tự động thành lớp và bất biến OCL trong cùng một hệ thống đối tượng
USE~\cite{gogolla2007use} dùng để kiểm i* và BPMN, nhờ đó ba mô hình được đối chiếu trên cùng một
trạng thái hệ thống duy nhất thay vì ba bản sao tách biệt.

Phương pháp đã được cài đặt thành một plugin tích hợp vào USE, gồm một bộ kiểm hàng loạt và một
trình gỡ lỗi trực quan phát lại từng bước trên hai diagram i*/BPMN đồng bộ. Công cụ được đánh giá
qua hai case study, một cuộc họp tổ chức nhiều vai trò và một quy trình ứng phó sự cố, cho thấy
phương pháp khả thi trên cả mô hình có lẫn không có ràng buộc vai trò.

Từ khoảng trống đã nêu hình thành hai câu hỏi nghiên cứu. \textbf{RQ1}: có phương pháp nào để đặc tả
cấu trúc hệ thống, gồm thực thể, vai trò, và ràng buộc giữa các vai trò, một cách hình thức ngay
trong giai đoạn phân tích, khi các ngôn ngữ mô hình ý định và mô hình tiến trình phổ biến đều không
có sẵn khái niệm này? \textbf{RQ2}: có phương pháp nào để kiểm tra tính nhất quán giữa yêu cầu, tức
mô hình ý định, và giải pháp, tức mô hình tiến trình, một cách tự động, trên cùng một trạng thái hệ
thống đã đặc tả cấu trúc tổ chức?

\textbf{Đóng góp.} Trong phạm vi cụ thể của cặp ngôn ngữ i*/BPMN và nền tảng USE/OCL, bài báo đóng
góp ba phần, ứng đúng thứ tự hai câu hỏi nghiên cứu trên.
\begin{enumerate}
  \item Trả lời RQ1: ngôn ngữ ACL đặc tả cấu trúc hệ thống, gồm thực thể, vai trò, nhóm, và ràng
    buộc tương thích vai trò, cùng luật dịch tự động sang mô hình lớp và bất biến OCL trong USE
    (Phần~\ref{sec:acl-language} và Phần~\ref{sec:acl2use});
  \item trả lời RQ2 trong phạm vi i*/BPMN: một thuật toán kiểm tra tính nhất quán yêu cầu--giải
    pháp dựa trên vết thực thi checkpoint, kế thừa khung kịch bản--snapshot nhưng suy snapshot tự
    động từ OCL, đồng thời tận dụng kết quả của đóng góp thứ nhất để snapshot phản ánh đúng cả
    ràng buộc cấu trúc (Phần~\ref{sec:checking});
  \item một công cụ hỗ trợ tích hợp cả hai đóng góp trên vào môi trường USE, gồm chế độ kiểm hàng
    loạt và một trình gỡ lỗi trực quan từng bước (Phần~\ref{sec:tool-support}).
\end{enumerate}

\textbf{Bố cục.} Phần~\ref{sec:related-work} khảo sát công trình liên quan theo năm nhóm cơ chế.
Phần~\ref{sec:motivation} đến Phần~\ref{sec:checking} trình bày phương pháp đề xuất, lần lượt qua
một ví dụ minh hoạ đầy đủ, kiến thức nền, tổng quan kiến trúc, chiến lược tiếp cận, ngôn ngữ ACL,
phép dịch ACL sang USE, và thuật toán kiểm tra. Phần~\ref{sec:discussion} thảo luận vị trí của
phương pháp so với các hướng đã khảo sát. Phần~\ref{sec:limitations} nêu hạn chế của phương pháp
đề xuất. Phần~\ref{sec:tool-support} trình bày công cụ hỗ trợ.
